\chapter{Éléments de théorie spectrale des graphes} \label{chap:graphs}

Cette annexe fournit des éléments de contexte pour le \Cref{theo:cheeger_inequality} établissant l'inégalité de Cheeger sur un graphe. Les notes des \Cref{sec:laplacian_of_a_graph,sec:eigenvalues_of_a_graph} proviennent de \cite[Sec. 1.2]{chung_lectures_1997} et celles de la \Cref{sec:the_cheeger_inequality} sont issues de \cite[Sec.~2]{ji_fourth_2010}.

\section{Laplacien d'un graphe} \label{sec:laplacian_of_a_graph}

Dans un graphe $\graphe$, notons $\degr(v)$ le degré du sommet $v$. Nous définissons le laplacien pour les graphes sans boucles ni arêtes multiples. Commençons par considérer la matrice $L$, définie comme suit :
\begin{equation}
L(u,v) \defeq \begin{cases}
    \degr(v) &\text{si } u = v, \\
    -1 &\text{si $u$ et $v$ sont adjacents}, \\
    0 & \text{sinon}.
\end{cases}
\end{equation}
Notons $T$ la matrice diagonale dont le coefficient d'indice $(v,v)$ vaut $\degr(v)$. Le \defini{laplacien} de $\graphe$ est défini par la matrice
\begin{equation}
\mathcal{L}(u,v) \defeq \begin{cases}
    1 & \text{si $u = v$ et $\degr(v) \not= 0$}, \\
    - \displaystyle\frac{1}{\sqrt{\degr(u) \degr(v)}} & \text{si $u$ et $v$ sont adjacents}, \\
    0 & \text{sinon}.
\end{cases}
\end{equation}
Nous pouvons écrire
\[
\mathcal{L} = T^{-1/2} \, L \, T^{-1/2}
\]
avec la convention $T^{-1}(v,v) = 0$ lorsque $\degr(v) = 0$. Le laplacien $\mathcal{L}$ peut être considéré comme un opérateur sur l'espace des fonctions $\fonctionens[g]{\vertices(\graphe)}{\R}$ vérifiant
\[
\mathcal{L}g(u) = \frac{1}{\sqrt{\degr(u)}} \sum\limits_{u \sim v} \bigg( \frac{g(u)}{\sqrt{\degr(u)}} - \frac{g(v)}{\sqrt{\degr(v)}} \bigg).
\]

\section{Valeurs propres d'un graphe} \label{sec:eigenvalues_of_a_graph}

Puisque $\mathcal{L}$ est symétrique, ses valeurs propres sont toutes réelles et positives. Nous pouvons utiliser les caractérisations variationnelles de ces valeurs propres à l'aide du quotient de Rayleigh de $\mathcal{L}$. Soit $g$ une fonction arbitraire qui associe à chaque sommet $v$ de $\graphe$ une valeur réelle $g(v)$. Nous pouvons considérer $g$ comme un vecteur colonne. Alors
\begin{align}
    \frac{\crodua{g}{\mathcal{L}g}}{\crodua{g}{g}} &= \frac{\crodua{g}{T^{-1/2} \, L \, T^{-1/2} g}}{\crodua{g}{g}}
    = \frac{\crodua{f}{Lf}}{\crodua{T^{1/2} f}{T^{1/2} f}}
    = \frac{\sum\limits_{u \sim v} \big( f(u) - f(v) \big)^2}{\sum\limits_v f(v)^2 \, \degr(v)} \label{eq:rayleigh_quotient}
\end{align}
où $g = T^{1/2} f$ et où $\sum_{u \sim v}$ désigne la somme sur toutes les paires non ordonnées $\{u, v\}$ telles que $u$ et $v$ sont adjacents. Ici, $\crodua{f}{g} = \sum_x f(x) \, g(x)$ désigne le produit scalaire usuel dans $\R^n$. La somme $\sum_{u \sim v} \big( f(u) - f(v) \big)^2$ est parfois appelée la \defini{somme de Dirichlet} de $\graphe$ et le rapport figurant au membre de gauche de \ref{eq:rayleigh_quotient} est souvent appelé le \defini{quotient de Rayleigh}.

L'\Cref{eq:rayleigh_quotient} montre que toutes les valeurs propres sont positives ou nulles. En fait, nous en déduisons aisément que $0$ est une valeur propre de $\mathcal{L}$. Notons les valeurs propres de $\mathcal{L}$ par $0 = \lambda_0 \leqslant \lambda_1 \leqslant \cdots \leqslant \lambda_{n-1}$. L'ensemble des $\lambda_i$ est appelé le \defini{spectre} de $\mathcal{L}$. Soit $\indicatrice$ la fonction constante qui vaut $1$ en chaque sommet. Alors $T^{1/2} \indicatrice$ est une fonction propre de $\mathcal{L}$ associée à la valeur propre $0$. De plus,
\begin{equation} \label{eq:lambdaG}
    \lambda_\graphe = \lambda_1 = \inf_{f \perp T \indicatrice} \frac{\sum\limits_{u \sim v} \big( f(u) - f(v) \big)^2}{\sum\limits_v f(v)^2 \, \degr(v)}.
\end{equation}
L'infimum est pris sur les fonctions $f$ définies sur les sommets de $\graphe$ et vérifiant $\sum\limits_v f(v) \, \degr(v) = 0$. La fonction propre correspondante est $g = T^{1/2} f$, comme dans l'\Cref{eq:rayleigh_quotient}. Le \defini{trou spectral} $\lambda_\graphe$ est la plus petite valeur propre non nulle de $\mathcal{L}$. Il est parfois commode de considérer la fonction non triviale $f$ réalisant l'infimum \ref{eq:lambdaG} ; nous appelons alors $f$ une \defini{fonction propre harmonique} de $\mathcal{L}$.

La formulation ci-dessus de $\lambda_\graphe$ correspond naturellement aux valeurs propres de l'opérateur de Laplace--Beltrami sur les variétés riemanniennes :
\[
\lambda_M = \inf_{\int_M f = 0} \frac{\int_M \abs{\grad f}^2}{\int_M \abs{f}^2}.
\]
Remarquons qu'ici la mesure associée à chaque arête vaut $1$. La mesure associée à chaque sommet est le degré de ce sommet. Notons que le trou spectral $\lambda_\graphe$ admet également la formulation
\[
\lambda_\graphe = \inf_f \frac{\sum\limits_{u \sim v} \big( f(u) - f(v) \big)^2}{\sum\limits_v \big( f(v) - \overline{f} \big)^2 \degr(v)}
\quad
\text{où}
\quad
\overline{f} = \displaystyle\frac{\sum\limits_v f(v) \degr(v)}{\vol(\graphe)}.
\]

\section{Inégalité de Cheeger} \label{sec:the_cheeger_inequality}

L'inégalité de Cheeger est l'un des outils majeurs de la théorie spectrale des graphes. Elle porte sur la relation entre deux invariants importants d'un graphe : le trou spectral et la constante de Cheeger. Pour un graphe non orienté $\graphe = (\vertices, \edges)$, nous notons $\lambda_\graphe$ le trou spectral du laplacien normalisé. Notons $h(\graphe)$ la constante de Cheeger de $\graphe$, définie comme la plus petite valeur $h$ telle que toute coupe isolant un ensemble de volume $x$ nécessite au moins $h \, x$ arêtes lorsque $x$ est inférieur à la moitié du volume total. L'inégalité de Cheeger affirme que, pour tout graphe connexe $\graphe$,
\[
2 \, h(\graphe) \geqslant \lambda_\graphe \geqslant \frac{h(\graphe)^2}{2}.
\]
La démonstration de l'inégalité de Cheeger (discrète) ci-dessus est très similaire à celle du cas continu en géométrie spectrale, due à Jeff Cheeger \cite{gunning_lower_2015}.

\begin{defi}[Constante de Cheeger]
Pour un ensemble de sommets $W \subset \vertices(\graphe)$, le \defini{rapport de Cheeger} de $W$ est défini par
\begin{equation}
h_W \defeq \frac{\abs{\partial W}}{\min\{ \vol(W), \vol(\graphe) - \vol(W) \}},
\end{equation}
où le bord de $W$ est noté $\partial W = \ens[\big]{ \{u, v\} \in \edges(\graphe) \tq u \in W \text{ et } v \not\in W }$. Notons $\comp{W}$ le complémentaire de $W$. Alors $h_W = h_{\comp{W}}$. La \defini{constante de Cheeger} de $\graphe$ est définie par
\begin{equation}
h(\graphe) \defeq \min_{W \subset \vertices(\graphe)} h_W.
\end{equation}
\end{defi}

\begin{remarque}
Remarquons que $\vol(\graphe) - \vol(W) = \vol(\comp{W})$.
\end{remarque}

\begin{theo}[Inégalité de Cheeger]
Dans un graphe $\graphe$, la constante de Cheeger $h(\graphe)$ et le trou spectral $\lambda_\graphe$ vérifient l'encadrement
\[
2 \, h(\graphe) \geqslant \lambda_\graphe \geqslant \frac{h(\graphe)^2}{2}.
\]
\end{theo}